\subsection{广群的正向系统}

在本节及下一节中，我们假设 I 是右有向的。

定义2. 关于指标集I的广群正向系统是一个集合的正向系统 \((\mathrm{E}_{\alpha}, f_{\beta\alpha})\) ，其中每个 \(E_{\alpha}\) 具有一个广群结构，且每个 \(f_{\beta\alpha}\) 都是广群同态。

设 \((\mathrm{E}_{\alpha}, f_{\beta\alpha})\) 是一个广群正向系统。用 E 表示正向极限集 \(\lim_{\rightarrow} \mathrm{E}_{\alpha}\) ，用 \(f_{\alpha}\) 表示从 \(\mathrm{E}_{\alpha}\) 到 E 的典范映射。回顾

\begin{enumerate}
\def\labelenumi{(\arabic{enumi})}
\setcounter{enumi}{1}
\tightlist
\item
\end{enumerate}

\[
f _ {\beta} \circ f _ {\beta \alpha} = f _ {\alpha} \quad \text { for } \alpha \leqslant \beta ,\tag{3}
\]

\[
\mathrm{E} = \bigcup_ {\alpha \in \mathbf {I}} f _ {\alpha} (\mathrm{E} _ {\alpha}).
\]

由(2)，同样地

\[
f _ {\alpha} (\mathrm{E} _ {\alpha}) \subset f _ {\beta} (\mathrm{E} _ {\beta}) \quad \text { for } \alpha \leqslant \beta .\tag{4}
\]

若 \(x_{\alpha}, y_{\alpha} \in \mathrm{E}_{\alpha}\) 满足 \(f_{\alpha}(x_{\alpha}) = f_{\alpha}(y_{\alpha})\) ，则存在 \(\beta \geqslant \alpha\) 使得 \(f_{\beta \alpha}(x_{\alpha}) = f_{\beta \alpha}(y_{\alpha})\) 。

命题1. 在 E 上存在一种且仅一种广群结构，使得映射 \(f_{\alpha} \colon \mathrm{E}_{\alpha} \to \mathrm{E}\) 是同态。若诸广群 \(\mathrm{E}_{\alpha}\) 是结合的（相应地，交换的），则 E 亦然。若诸广群 \(\mathrm{E}_{\alpha}\) 与诸同态 \(f_{\beta \alpha}\) 是保单位的，则广群 E 与诸同态 \(f_{\alpha}\) 亦然。

诸广群 \(E_{\alpha}\) 将以乘法记法书写。

设 \(x, y\) 是 E 中的元素。存在属于 I 的 \(\alpha\) 以及 \(x_{\alpha}, y_{\alpha}\) 属于 \(E_{\alpha}\) ，使得 \(x = f_{\alpha}(x_{\alpha})\) 与 \(y = f_{\alpha}(y_{\alpha})\) 。若 E 上存在使 \(f_{\alpha}\) 为同态的广群结构，则 \(x.y = f_{\alpha}(x_{\alpha}y_{\alpha})\) ，由此得该广群结构的唯一性。

为证明存在性，我们必须证明：对 I 中的 \(\alpha, \beta\) ， \(x_{\alpha}, y_{\alpha}\) 属于 E \(_{\alpha}\) 且 \(x_{\beta}', y_{\beta}'\) 属于 E \(_{\beta}\) 时，关系

\[
f _ {\alpha} (x _ {\alpha}) = f _ {\beta} (x _ {\beta} ^ {\prime}), \quad f _ {\alpha} (y _ {\alpha}) = f _ {\beta} (y _ {\beta} ^ {\prime})\tag{5}
\]

蕴含 \(f_{\alpha}(x_{\alpha}y_{\alpha}) = f_{\beta}(x_{\beta}'y_{\beta}')\) 。对于 \(\gamma \geqslant \alpha\) 和 \(\gamma \geqslant \beta\) ，我们设 \(x_{\gamma} = f_{\gamma \alpha}(x_{\alpha})\) , \(y_{\gamma} = f_{\gamma \alpha}(y_{\alpha}), x_{\gamma}^{\prime} = f_{\gamma \beta}(x_{\beta}^{\prime}), y_{\gamma}^{\prime} = f_{\gamma \beta}(y_{\beta}^{\prime})\) 。根据正向极限的定义，存在 I 中的 \(\gamma\) 满足 \(\gamma \geqslant \alpha, \gamma \geqslant \beta, x_{\gamma} = x_{\gamma}^{\prime}, y_{\gamma} = y_{\gamma}^{\prime}\) 。则

\[
\begin{array}{r l} f _ {\alpha} (x _ {\alpha} y _ {\alpha}) & = f _ {\gamma} (f _ {\gamma \alpha} (x _ {\alpha} y _ {\alpha})) = f _ {\gamma} (x _ {\gamma} y _ {\gamma}) = f _ {\gamma} (x _ {\gamma} ^ {\prime} y _ {\gamma} ^ {\prime}) = f _ {\gamma} (f _ {\gamma \beta} (x _ {\beta} ^ {\prime} y _ {\beta} ^ {\prime})) \\ & = f _ {\beta} (x _ {\beta} ^ {\prime} y _ {\beta} ^ {\prime}). \end{array}
\]

假设诸广群 \(E_{\alpha}\) 是结合的。设 x、y、z 属于 E。存在 \(\alpha \in I\) 及 \(x_{\alpha}, y_{\alpha}, z_{\alpha}\) 属于 \(E_{\alpha}\) ，使得

\[
x = f _ {\alpha} (x _ {\alpha}), \quad y = f _ {\alpha} (y _ {\alpha}), \quad z = f _ {\alpha} (z _ {\alpha}).
\]

则 \(xy = f_{\alpha}(x_{\alpha}y_{\alpha})\) ，从而 \((xy)z = f_{\alpha}((x_{\alpha}y_{\alpha})z_{\alpha})\) ；类似地

\[
x (y z) = f _ {\alpha} \left(x _ {\alpha} \left(y _ {\alpha} z _ {\alpha}\right)\right),
\]

从而 \((xy)z = x(yz)\) ，因为 \((x_{\alpha}y_{\alpha})z_{\alpha} = x_{\alpha}(y_{\alpha}z_{\alpha})\) 。交换广群的情形可类似地处理。

最后假设每个广群 \(E_{\alpha}\) 有单位元 \(e_{\alpha}\) ，且 \(f_{\beta\alpha}(e_{\alpha}) = e_{\beta}\) 对 \(\alpha \leqslant \beta\) 成立。对 I 中的 \(\alpha, \beta\) ，存在 I 中的 \(\gamma\) 满足 \(\gamma \geqslant \alpha\) 与 \(\gamma \geqslant \beta\) ，从而

\[
f _ {\alpha} (e _ {\alpha}) = f _ {\gamma} \left(f _ {\gamma \alpha} (e _ {\alpha})\right) = f _ {\gamma} (e _ {\gamma}) = f _ {\gamma} \left(f _ {\gamma \beta} (e _ {\beta})\right) = f _ {\beta} (e _ {\beta})
\]

因此存在元素 \(e\) 于 \(\mathbf{E}\) 中，使得 \(f_{\alpha}(e_{\alpha}) = e\) 对一切 \(\alpha \in \mathbf{I}\) 成立。设 \(x\) 是 \(\mathbf{E}\) 中的元素；选取 \(\alpha \in \mathbf{I}\) 与 \(x_{\alpha} \in \mathrm{E}_{\alpha}\) ，使得 \(x = f_{\alpha}(x_{\alpha})\) 。则

\[
e x = f _ {\alpha} (e _ {\alpha}) \cdot f _ {\alpha} (x _ {\alpha}) = f _ {\alpha} (e _ {\alpha}. x _ {\alpha}) = f _ {\alpha} (x _ {\alpha}) = x
\]

并且类似地有 \(x.e = x\) ，因此 \(e\) 是 \(E\) 的单位元。

广群 E 称为诸广群 \(E_{\alpha}\) 的正向极限。

命题2. 设 \((\mathrm{E}_{\alpha}, f_{\beta \alpha})\) 是一个广群正向系统，E 为其正向极限， \(f_{\alpha} \colon \mathrm{E}_{\alpha} \to \mathrm{E}\) 为典范同态。假设给定一个广群 F 及同态 \(u_{\alpha} \colon \mathrm{E}_{\alpha} \to \mathrm{F}\) ，使得 \(u_{\alpha} = u_{\beta} \circ f_{\beta \alpha}\) 对 \(\alpha \leqslant \beta\) 成立。存在一个且仅一个同态 \(u \colon \mathrm{E} \to \mathrm{F}\) ，使得 \(u_{\alpha} = u \circ f_{\alpha}\) 对所有 \(\alpha \in \mathbf{I}\) 成立。若诸广群 \(\mathrm{E}_{\alpha}\) 与 F 以及诸同态 \(f_{\beta \alpha}\) 与 \(u_{\alpha}\) 都是保单位的，则同态 \(u\) 是保单位的。

我们知道（《集合论》III，§ 7，第 6 号，命题6）存在一个且仅一个映射 \(u: \mathrm{E} \to \mathrm{F}\) ，使得 \(u_{\alpha} = u \circ f_{\alpha}\) 对所有 \(\alpha \in \mathbf{I}\) 成立。我们验证 \(u\) 是一个同态：设 \(x, y\) 属于 \(\mathrm{E}, \alpha\) 属于 \(\mathrm{I}\) ，且 \(x_{\alpha}, y_{\alpha}\) 属于 \(\mathrm{E}_{\alpha}\) ，满足 \(x = f_{\alpha}(x_{\alpha})\) 与 \(y = f_{\alpha}(y_{\alpha})\) 。则 \(xy = f_{\alpha}(x_{\alpha}y_{\alpha})\) ，从而

\[
\begin{array}{r l} u (x y) & = u (f _ {\alpha} (x _ {\alpha} y _ {\alpha})) = u _ {\alpha} (x _ {\alpha} y _ {\alpha}) = u _ {\alpha} (x _ {\alpha}) u _ {\alpha} (y _ {\alpha}) \\ & = u (f _ {\alpha} (u _ {\alpha})) u (f _ {\alpha} (y _ {\alpha})) = u (x) u (y). \end{array}
\]

现在考虑保单位的情形，并用 \(e_{\alpha}\) 表示 \(E_{\alpha}\) 的单位元， \(e\) 表示 \(E\) 的单位元， \(e'\) 表示 \(F\) 的单位元。选取 \(\alpha \in I\) ，则 \(e = f_{\alpha}(e_{\alpha})\) ，从而

\[
u (e) = u \left(f _ {\alpha} \left(e _ {\alpha}\right)\right) = u _ {\alpha} \left(e _ {\alpha}\right) = e ^ {\prime}
\]

因为 \(u_{\alpha}\) 是保单位的。因此 \(u\) 是保单位的。

与广群正向系统的概念类似，可以表述幺半群或群的正向系统的概念。命题1表明，作为幺半群正向系统 \((\mathrm{E}_{\alpha}, f_{\beta \alpha})_{\alpha, \beta \in I}\) 之极限的广群 E 是一个幺半群。我们证明，若诸 \(\mathrm{E}_{\alpha}\) 是群，则 E 是群：设 \(x \in \mathrm{E}\) , \(\alpha \in I\) 与 \(x_{\alpha} \in \mathrm{E}_{\alpha}\) 满足 \(x = f_{\alpha}(x_{\alpha})\) ；E 的元素 \(y = f_{\alpha}(x_{\alpha}^{-1})\) 是 \(x\) 的逆元（§ 2，第 3 号）。命题2的泛性质立即推广到幺半群或群的正向系统的情形。

留给读者定义环的正向系统。设 \((\mathbf{A}_{\alpha}, f_{\beta \alpha})\) 是这样一个正向系统；令 \(\mathbf{A} = \lim_{\longrightarrow} \mathbf{A}_{\alpha}\) ，并令 \(f_{\alpha}: \mathbf{A}_{\alpha} \to \mathbf{A}\) 为典范同态。在 \(\mathbf{A}\) 上存在（命题2）一种且仅一种加法与乘法，由 \(x + y = f_{\alpha}(x_{\alpha} + y_{\alpha})\) , \(xy = f_{\alpha}(x_{\alpha}y_{\alpha})\) （对 \(\alpha \in \mathbf{I}\) 、 \(x_{\alpha}, y_{\alpha}\) 属于 \(\mathbf{A}_{\alpha}\) 且 \(x = f_{\alpha}(x_{\alpha})\) , \(y = f_{\alpha}(y_{\alpha})\) ）刻画。在加法下 \(\mathbf{A}\) 是交换群，乘法是结合的且保单位的。最后，对 \(x, y, z\) 属于 \(\mathbf{A}\) ，选取 \(\alpha\) 属于 \(\mathbf{I}\) 及 \(x_{\alpha}, y_{\alpha}\) 与 \(z_{\alpha}\) 属于 \(\mathbf{A}_{\alpha}\) ，使得

\[
x = f _ {\alpha} (x _ {\alpha}), \quad y = f _ {\alpha} (y _ {\alpha}) \quad \text { and } \quad z = f _ {\alpha} (z _ {\alpha}).
\]

那么

\[
\begin{array}{r l} (x + y) \cdot z & = f _ {\alpha} (x _ {\alpha} + y _ {\alpha}) f _ {\alpha} (z _ {\alpha}) = f _ {\alpha} ((x _ {\alpha} + y _ {\alpha}) z _ {\alpha}) \\ & = f _ {\alpha} (x _ {\alpha} z _ {\alpha} + y _ {\alpha} z _ {\alpha}) = f _ {\alpha} (x _ {\alpha} z _ {\alpha}) + f _ {\alpha} (y _ {\alpha} z _ {\alpha}) = x z + y z \end{array}
\]

而关系 \(x(y + z) = xy + xz\) 可类似地证明。换言之，A 具有一个环结构，其特征由 \(f_{\alpha}\) 对所有 \(\alpha \in I\) 是环同态这一事实刻画。

环 A 称为诸环 \(A_{\alpha}\) 的正向极限。命题2立即推广到环的情形。

命题3. (a) 若诸 \(\mathbf{A}_{\alpha}\) 非零，则 \(\mathbf{A}\) 非零。

\begin{enumerate}
\def\labelenumi{(\alph{enumi})}
\setcounter{enumi}{1}
\item
  若诸 \(\mathbf{A}_{\alpha}\) 都是整环，则 \(\mathbf{A}\) 是整环。
\item
  若诸 \(\mathbf{A}_{\alpha}\) 都是域，则 \(\mathbf{A}\) 是域。
\end{enumerate}

设 \(0_{\alpha}, 1_{\alpha}\) 为 \(A_{\alpha}\) 的零元与单位元，0、1 为 A 的零元与单位元。存在 \(\alpha \in I\) 使得 \(f_{\alpha}(0_{\alpha}) = 0, f_{\alpha}(1_{\alpha}) = 1\) 。若 0 = 1，则存在 \(\beta \geqslant \alpha\) 使得 \(f_{\beta \alpha}(0_{\alpha}) = f_{\beta \alpha}(1_{\alpha})\) ，即 \(0_{\beta} = 1_{\beta}\) 。这就证明了 (a)。

假设诸 \(A_{\alpha}\) 是整环。则由 (a)，A 是交换的且非零。设 x、y 为 A 中满足 xy = 0 的元素。存在 \(\alpha \in I\) 与 \(x_{\alpha}, y_{\alpha} \in A_{\alpha}\) ，使得 \(x = f_{\alpha}(x_{\alpha}), y = f_{\alpha}(y_{\alpha})\) 。则 \(f_{\alpha}(x_{\alpha}y_{\alpha}) = xy = 0 = f_{\alpha}(0_{\alpha})\) 。因此存在 \(\beta \geqslant \alpha\) 使得 \(f_{\beta\alpha}(x_{\alpha}y_{\alpha}) = f_{\beta\alpha}(0_{\alpha})\) 。由于 \(A_{\beta}\) 是整环，由此得 \(f_{\beta\alpha}(x_{\alpha}) = 0_{\beta}\) 或 \(f_{\beta\alpha}(y_{\alpha}) = 0_{\beta}\) ，从而 x = 0 或 y = 0。这就证明了 (b)。

假设诸 \(A_{\alpha}\) 是域。则由 (a)， \(A \neq \{0\}\) 。设 x 是 A 的一个非零元素。存在 \(\alpha \in I\) 与 \(x_{\alpha} \in A_{\alpha}\) ，使得 \(x = f_{\alpha}(x_{\alpha})\) 。则 \(x_{\alpha} \neq 0\) 且 \(f_{\alpha}(x_{\alpha}^{-1})\) 是 x 在 A 中的逆元。这就证明了 (c)。

设 \(\mathfrak{E} = (\mathrm{E}_{\alpha}, f_{\beta \alpha})\) 与 \(\mathfrak{E}' = (\mathrm{E}_{\alpha}', f_{\beta \alpha}')\) 是广群（相应地，幺半群、群、环）的两个正向系统。从 \(\mathfrak{E}\) 到 \(\mathfrak{E}'\) 的一个同态是一个映射的正向系统 \((u_{\alpha})_{\alpha \in I}\) （ \(u_{\alpha} \colon E_{\alpha} \to E_{\alpha}'\) ），使得每个 \(u_{\alpha}\) 都是同态。在这些条件下，映射 \(u = \lim u_{\alpha}\) 从 \(E = \lim E_{\alpha}\) 到 \(E' = \lim E_{\alpha}'\) 是一个同态（参见《集合论》III，§ 7，第 6 号）。

